\section{Conclusion and Future Work}
\label{sec:conclusion}

In this paper, we introduced Curvature-Aware Analytical Model Merging (ACM), a mathematically rigorous and training-free framework for multi-task parameter consolidation. Moving beyond empirical heuristics like test-time unsupervised entropy minimization, ACM formulates model merging as a quadratic minimization problem derived from second-order Taylor expansions of task losses. By projecting the high-dimensional weight space onto the low-dimensional task vector subspace, we successfully capture the full, non-diagonal, cross-parameter Hessian curvature with zero diagonal approximation. Under Ridge regularization, this formulation yields an exact closed-form analytical solution solved in a single step with zero test-time training overhead.

Extensive simulation sweeps across 30 seeds show that in highly realistic, coupled non-convex environments, ACM outperforms existing test-time adaptation methods by up to +8.11\%. Physical validation on Vision Transformers confirms that our proposed leakage-free Global-Normalized ACM (ACM-GlobalNorm) achieves a Joint Average accuracy of 57.76\%, significantly outperforming diagonal Fisher Merging (56.03\%) and polynomial test-time adaptation (38.96\%) without any test-time optimization. When compared with an exhaustively tuned Task Arithmetic baseline (which achieves 60.72\% at scale 0.4), we analyze the limitations of the local quadratic Taylor approximation on highly non-convex, converged physical neural manifolds. We highlight that as expert task vectors grow in magnitude during convergence, the approximation error of local curvature-based methods increases, explaining this local-global gap. This provides a deep, mathematically transparent foundation for future work, emphasizing the necessity of exploring higher-order expansions or dynamic regularized calibration to bridge the local-global optimization gap. Finally, our layer-wise coefficient analysis demonstrates ACM's deep physical compliance, including the automatic detection of untrained layers and the mathematical derivation of negative scaling factors that actively cancel out destructive representation interference.

\subsection{Future Directions}
While ACM provides a powerful and mathematically robust foundation, several promising future directions exist:
\begin{enumerate}
    \item \textbf{Higher-Order Taylor Approximations:} Incorporating third-order or fourth-order tensor derivatives to better capture extreme landscape non-convexities, solved via tensor decomposition or zero-order adjustments.
    \item \textbf{Dynamic Calibration Streams:} Extending the analytical system to support streaming target distributions where the task priors and scaling coefficients are dynamically adjusted online based on real-time task vector activation signals.
    \item \textbf{Cross-Layer Hessian Interactions:} Modeling block-tridiagonal or block-sparse cross-layer Hessian correlations using hierarchical solvers, further relaxing Assumption 3.1.
\end{enumerate}
In conclusion, we believe that bringing mathematical rigor and formal geometric guarantees to parameter fusion is essential for the long-term success of the model merging paradigm. We hope that ACM inspires the community to explore deeper theoretical avenues in deep learning weight consolidation.
